\documentclass[../../main.tex]{subfiles}
\begin{document}
\begin{oframed}
  \subfile{problem}
\end{oframed}

{\bf [解答]}

赤札R, 白札Wと表す．人が$W$，$R$を選ぶ確率は確率$\dfrac{1}{2}$で同様に確からしい．

\paragraph{(1)}

あり得る状態を全てリストアップすると以下のようになる．

\begin{table}[htb]
  \centering
  \caption{札の個数による状態の分類}
  \begin{tabular}{|c|c|c|c|l|}
    \hline
    & A & B & C & \\ \hline
    (R,W) & (1,1) & (1,1) & (1,1) & $F_0$ \\ \hline
    (R,W) & (1,1) & (0,2) & (2,0) & $F_1$ \\ \hline
    (R,W) & (2,0) & (0,2) & (1,1) & $F_2$ \\ \hline
    (R,W) & (1,1) & (2,0) & (0,2) & $F_3$ \\ \hline
    (R,W) & (0,2) & (2,0) & (1,1) & $F_4$ \\ \hline
    (R,W) & (0,2) & (1,1) & (2,0) & $F_5$ \\ \hline
    (R,W) & (2,0) & (1,1) & (0,2) & $F_6$ \\ \hline
  \end{tabular}
\end{table}

\paragraph{(2)}

遷移確率をリストアップすると以下のようになる．
\begin{align}
  \begin{cases}
    P_{00} = \dfrac{1}{4} \\
    P_{0k} = \dfrac{1}{8} \ (k=1,2,\cdots,6)
  \end{cases}
\end{align}
\begin{align}
  \begin{cases}
    P_{1k} = 0 \quad (k=1,3,4,5,6) \\
    P_{10} = P_{12} = \dfrac{1}{2}
  \end{cases}
\end{align}
\begin{align}
  \begin{cases}
    P_{2k} = 0 \ (k=1,2,3,4,5) \\
    P_{20} = P_{26} = \dfrac{1}{2}
  \end{cases}
\end{align}
\begin{align}
  \begin{cases}
    P_{3k} = 0 \ (k=1,2,3,5,6) \\
    P_{30} = P_{34} = \dfrac{1}{2}
  \end{cases}
\end{align}
\begin{align}
  \begin{cases}
    P_{4k} = 0 \ (k=1,3,4,5,6) \\
    P_{40} = P_{45} = \dfrac{1}{2}
  \end{cases}
\end{align}
\begin{align}
  \begin{cases}
    P_{5k} = 0 \ (k=2,3,4,5,6) \\
    P_{50} = P_{51} = \dfrac{1}{2}
  \end{cases}
\end{align}
\begin{align}
  \begin{cases}
    P_{6k} = 0 \ (k=1,2,4,5,6) \\
    P_{60} = P_{63} = \dfrac{1}{2}
  \end{cases}
\end{align}

\paragraph{(3)}

\begin{figure}[htb]
  \centering
  \begin{tikzpicture}[x=4.8cm, y=1.6cm, >=stealth, every node/.style={font=\small}]

    % --- 1. 列ヘッダー ---
    \node[font=\normalsize\bfseries, gray!30!black] at (0, 1.3) {$n$ 回目};
    \node[font=\normalsize\bfseries, gray!30!black] at (1, 1.3) {$n+1$ 回目};

    % --- 2. 状態ノード ---
    % n 回目
    \node[draw=blue!85!black, fill=blue!6, thick, rounded corners=4pt, inner sep=6pt, align=center] (F0_n) at (0, 0.5) {
      \textbf{状態 $F_0$}\\
      \footnotesize (全員が赤1枚・白1枚)\\[1pt]
      \footnotesize 確率 $q_n$
    };

    \node[draw=teal!85!black, fill=teal!6, thick, rounded corners=4pt, inner sep=6pt, align=center] (Fother_n) at (0, -0.9) {
      \textbf{$F_0$ 以外の状態 ($F_1 \sim F_6$)}\\
      \footnotesize 確率 $1 - q_n$
    };

    % n+1 回目
    \node[draw=blue!85!black, fill=blue!6, thick, rounded corners=4pt, inner sep=6pt, align=center] (F0_next) at (1, 0.5) {
      \textbf{状態 $F_0$}\\
      \footnotesize 確率 $q_{n+1}$
    };

    \node[draw=teal!85!black, fill=teal!6, thick, rounded corners=4pt, inner sep=6pt, align=center] (Fother_next) at (1, -0.9) {
      \textbf{$F_0$ 以外の状態}\\
      \footnotesize 確率 $1 - q_{n+1}$
    };

    % --- 3. 遷移矢印と確率 ---
    % F0 -> F0
    \draw[->, line width=1.4pt, blue!85!black] (F0_n.east) -- (F0_next.west)
    node[midway, above=2pt, font=\normalsize\bfseries] {$\frac{1}{4}$};

    % F0 -> F_other
    \draw[->, line width=1.1pt, teal!80!black] (F0_n.east) -- (Fother_next.west)
    node[pos=0.7, below left=1pt, font=\normalsize\bfseries] {$\frac{3}{4}$};

    % F_other -> F0
    \draw[->, line width=1.4pt, red!85!black] (Fother_n.east) -- (F0_next.west)
    node[pos=0.7, above left=1pt, font=\normalsize\bfseries] {$\frac{1}{2}$};

    % F_other -> F_other
    \draw[->, line width=1.1pt, gray!70!black] (Fother_n.east) -- (Fother_next.west)
    node[midway, below=2pt, font=\normalsize\bfseries] {$\frac{1}{2}$};

    % --- 4. 注釈 ---
    \node[right=6pt, font=\footnotesize, blue!90!black] at (F0_next.east)
    {$\left( p_{00} = \frac{1}{4} \right)$};
    \node[right=6pt, font=\footnotesize, red!90!black] at (Fother_next.east)
    {$\left( p_{k0} = \frac{1}{2} \ (k=1,\dots,6) \right)$};

  \end{tikzpicture}
  \caption{$n$ 回目から $n+1$ 回目への状態遷移図}
  \label{fig:state_transition_1987_6}
\end{figure}

$q_n$の漸化式について考える．(2)で
\begin{align}
  P_{00} = \dfrac{1}{4}, P_{k0} = \dfrac{1}{2} \quad (k=1,2,\cdots,6)
\end{align}
だから，$n$回目に$F_{0}$だったものが$n+1$回目で$F_{0}$に遷移する確率は$\dfrac{1}{4}$，$n$回目に$F_{0}$以外だったものが$n+1$回目で$F_{0}$に遷移する確率は$\dfrac{1}{2}$だから，漸化式は
\begin{align}
  \begin{cases}
    q_{n+1} = \dfrac{1}{4}q_n + \dfrac{1}{2}(1-q_n)
    = -\dfrac{1}{4}q_n + \dfrac{1}{2} \\
    q_0 = 1
  \end{cases}
\end{align}
となる．等比数列の公式から
\begin{align}
  q_n
  &= \left(-\dfrac{1}{4}\right)^n \left(1 - \dfrac{2}{5}\right) + \dfrac{2}{5} \\
  &= \dfrac{3}{5}\left(-\dfrac{1}{4}\right)^n + \dfrac{2}{5}
\end{align}
となる．

\end{document}
